\documentclass[12pt,a4paper]{article}

% === Packages ===
\usepackage[margin=2.5cm]{geometry}
\usepackage{amsmath,amssymb,amsfonts}
\usepackage{booktabs}
\usepackage{graphicx}
\usepackage{hyperref}
\usepackage{setspace}
\usepackage{titlesec}
\usepackage{url}
\usepackage{array}
\usepackage{float}
\usepackage{caption}

% === Line Spacing ===
\onehalfspacing

% === Hyperref Setup ===
\hypersetup{
    colorlinks=true,
    linkcolor=blue,
    citecolor=blue,
    urlcolor=blue,
    pdftitle={Beyond the Operations Center: A Unified Co-evolutionary Theory of NOC-to-SOC Migration},
    pdfauthor={Anonymous}
}

% === Section Formatting ===
\titleformat{\section}{\normalfont\Large\bfseries}{\thesection}{1em}{}
\titleformat{\subsection}{\normalfont\large\bfseries}{\thesubsection}{1em}{}
\titleformat{\subsubsection}{\normalfont\normalsize\bfseries}{\thesubsubsection}{1em}{}

% === Page Numbers ===
\pagestyle{plain}
\pagenumbering{arabic}

\begin{document}

% ===================================================================
% TITLE PAGE
% ===================================================================
\begin{center}
\vspace*{2cm}

{\LARGE\bfseries Beyond the Operations Center:\\[0.5cm]
A Unified Co-evolutionary Theory of NOC-to-SOC Migration\\in the Telecommunications Sector}

\vspace{1.5cm}

{\large\itshape The COMET Framework: Co-evolutionary Operations Maturity\\and Epistemic Transformation}

\vspace{2cm}

{\normalsize Research Paper\\[0.3cm]
Technology Strategy \& Operations Management\\[0.3cm]
June 2026}

\vspace{3cm}

{\small\textit{This paper synthesizes and extends three prior theoretical contributions---ISDAF, CMOR, and SCOTT---into a unified framework with formal mathematical models and testable hypotheses for NOC-to-SOC migration in telecommunications.}}

\end{center}

\newpage

% ===================================================================
% TABLE OF CONTENTS
% ===================================================================
\tableofcontents
\newpage

% ===================================================================
% ABSTRACT
% ===================================================================
\section{Abstract}

The migration from Network Operations Centers (NOCs) to Service Operations Centers (SOCs) constitutes one of the most consequential structural transformations in the global telecommunications industry. Despite the proliferation of industry frameworks---most notably the TM Forum's Autonomous Operations Maturity Model (AOMM), Value Operations Framework (VOF), and the Measuring and Managing Autonomy (MAMA) toolkit---theoretical explanations of this migration remain fragmented across three competing perspectives: the sociotechnical integration perspective (ISDAF), the contextual maturity perspective (CMOR), and the dialectical co-evolution perspective (SCOTT). Each framework captures important but partial dimensions of the transformation, and critical contradictions between them remain unresolved. This paper addresses these deficiencies by developing the \textbf{Co-evolutionary Operations Maturity and Epistemic Transformation (COMET)} theory, a unified framework that synthesizes the strongest elements of all three prior contributions while resolving their theoretical tensions and extending them through formal mathematical modeling. COMET is grounded in three foundational pillars---Dynamic Capabilities Theory, Socio-Technical Systems Theory, and Boundary Object Theory---and organized around seven interconnected dimensions: Technological Maturity, Epistemic Infrastructure, Data Fabric Maturity, Cultural Readiness, Organizational Morphology, Contextual Maturity, and Value Governance. The paper contributes ten formal mathematical models, including formulations for Expected Demand Not Served (EDNS), Service Graph Maturity, Automation Velocity, Maturity Coherence Index, Cultural Readiness Composite, a contextual interaction regression model, a Data Fabric--MTTR exponential decay function, a cost-benefit model for closed-loop automation, and a Markov transition model for maturity progression. Eight testable hypotheses are derived, addressing the relative predictive power of Maturity Coherence versus Cultural Readiness, the mediating role of Service Graph topology, the moderating effects of contextual factors, and the non-linear dynamics of transformation trajectories. The COMET framework advances the theoretical discourse by demonstrating that NOC-to-SOC migration is neither a linear technological substitution nor a simple co-evolutionary process, but a structurally coupled system in which epistemic, organizational, cultural, and contextual dimensions co-determine transformation outcomes through mathematically specifiable mechanisms.

\vspace{0.3cm}
\noindent\textbf{Keywords:} Service Operations Center (SOC), Network Operations Center (NOC), Autonomous Networks, AOMM, Digital Twin Network, Data Fabric, Closed-Loop Automation, TM Forum, Dynamic Capabilities, Socio-Technical Systems, Maturity Models, Telecommunications Operations

\newpage

% ===================================================================
% 1. INTRODUCTION
% ===================================================================
\section{Introduction}

The telecommunications industry stands at a critical inflection point. The exponential growth in data traffic, coupled with the proliferation of advanced digital services such as 5G network slicing, software-defined wide-area networking (SD-WAN), and edge computing, has fundamentally altered the relationship between network infrastructure and customer experience. As data traffic grows exponentially while revenue remains flat---a phenomenon termed the ``traffic-revenue decoupling''---Communications Service Providers (CSPs) are compelled to reconceptualize their operational models \citep{tmforum2024a}. The emerging paradigm is the Service Operations Center (SOC), a construct defined by the TM Forum as ``a central hub within a telecom operator's Customer Experience Management (CEM) strategy'' that ``oversees, monitors, and manages the delivery of services in an organization'' \citep[p.~3]{tmforum2025soc}.

At its core, the NOC-to-SOC migration represents a specific manifestation of a broader tension in technology-intensive industries: the persistent gap between technical systems and the human experiences those systems are designed to serve. The NOC, by its very design, operates within the logic of resource management, optimizing for network element health, infrastructure availability, and technical key performance indicators (KPIs) such as packet loss, latency, and throughput. The SOC, conversely, operates within the logic of experience management, optimizing for service quality as perceived by end users, business outcomes such as churn reduction and revenue protection, and customer-centric metrics such as the Key Quality Indicator (KQI), Customer Experience Index (CEI), and Net Promoter Score (NPS). The migration from one to the other thus constitutes not merely a technological upgrade but an epistemological reorientation---a shift in what counts as knowledge, what counts as value, and what counts as success.

The TM Forum has been the primary institutional driver of this transformation, developing an ecosystem of interconnected frameworks that provide both the conceptual vocabulary and the implementation guidance for SOC adoption. The Autonomous Operations Maturity Model (AOMM) \citep{tmforum2024b} defines a six-level progression from Level~0 (manual operations) through Level~5 (fully autonomous), assessing maturity across seven operational domains: Fault Management, Performance Management, Configuration Management, Change Management, Security Management, Service Fulfillment, and Service Assurance. The Expected Demand Not Served (EDNS) metric \citep{tmforum2024c} reframes the traditional availability paradigm by measuring service impact in terms of demand-weighted quality degradation rather than binary up/down states. The Measuring and Managing Autonomy (MAMA) framework \citep{tmforum2024d} provides operational guidance for autonomy deployment, while the R.I.S.E.\ value model (Revenue growth, Innovation acceleration, Service experience enhancement, and Efficiency optimization) articulated in the primary strategic data \citep{primary2026a,primary2026b} connects operational transformation to business value creation.

Despite the strategic importance of this transformation, the academic literature has been conspicuously fragmented. Three recent theoretical contributions have emerged, each advancing a distinct perspective on the migration. The Integrated Sociotechnical-Dynamic Autonomy Framework (ISDAF) emphasizes the joint optimization of technical and social subsystems through the lens of Sociotechnical Systems Theory and Dynamic Capabilities Theory. The Contextual Maturity and Organizational Readiness (CMOR) framework extends the AOMM by incorporating organizational culture, human capital readiness, and market contextual factors. The Service-Centric Operations Transformation Theory (SCOTT) introduces the concepts of epistemic infrastructure, morphological elasticity, process genealogy, and boundary object mediation. While each framework advances the theoretical discourse, they exhibit notable gaps: ISDAF lacks testable hypotheses and formal mathematical models; CMOR contains no quantitative specifications despite introducing six hypotheses; and SCOTT, while theoretically the most sophisticated, provides no mathematical formalization of its key constructs. Furthermore, contradictions between the frameworks---regarding the relative importance of cultural readiness versus maturity coherence, the linearity versus non-linearity of migration trajectories, and the treatment of the NOC-SOC relationship as replacement versus co-evolution---remain unresolved.

This paper addresses these theoretical deficiencies by developing the Co-evolutionary Operations Maturity and Epistemic Transformation (COMET) theory. COMET synthesizes the strongest elements of ISDAF, CMOR, and SCOTT into a unified framework, resolves their principal contradictions through the concept of structural coupling, and extends the theoretical discourse through ten formal mathematical models and eight testable hypotheses. The paper proceeds as follows: Section~2 reviews the theoretical foundations; Section~3 presents the critical synthesis of the three prior frameworks; Section~4 develops the COMET theory; Section~5 presents the mathematical models; Section~6 discusses implications; and Section~7 concludes.

% ===================================================================
% 2. THEORETICAL FOUNDATIONS
% ===================================================================
\section{Theoretical Foundations}

\subsection{Dynamic Capabilities Theory}

Dynamic Capabilities Theory (DCT), as articulated by \citet{teece1997dynamic} and subsequently refined by \citet{teece2007explicating} and \citet{teece2018business}, provides the primary theoretical lens for understanding how telecommunications organizations can purposefully create, extend, or modify their resource base in response to rapidly changing environments. The theory distinguishes between ordinary capabilities, which permit an organization to earn a living in a given environment, and dynamic capabilities, which enable it to adapt, integrate, and reconfigure internal and external competences to address rapidly changing environments \citep{teece1997dynamic}. The three foundational micro-foundations---sensing (the ability to scan and interpret the environment), seizing (the ability to mobilize resources to capture value from opportunities), and reconfiguring (the ability to maintain evolutionary fitness by adapting the organizational architecture)---map directly onto the NOC-to-SOC transformation. The sensing dimension corresponds to the organization's capacity to recognize the limitations of network-centric operational models and to identify the service-centric opportunity. The seizing dimension corresponds to the investment decisions and organizational redesign efforts required to build SOC capabilities. The reconfiguring dimension corresponds to the ongoing adaptation of the operational model as technologies, markets, and customer expectations evolve.

The application of DCT to the NOC-to-SOC context is further enriched by the concept of absorptive capacity \citep{cohen1990absorptive}, which explains the organizational prerequisites for recognizing the value of new information, assimilating it, and applying it to commercial ends. In the SOC transformation context, absorptive capacity determines an operator's ability to adopt and effectively utilize advanced technologies such as Data Fabric, Digital Twin Networks, and generative AI. Operators with higher absorptive capacity---typically those with stronger R\&D functions, more diverse knowledge networks, and greater prior experience with digital transformation---are expected to navigate the SOC transition more effectively than those with lower capacity.

\subsection{Socio-Technical Systems Theory}

Socio-Technical Systems (STS) Theory, originating in the seminal work of \citet{trist1951social} and subsequently developed by \citet{bostrom1977mis} and others, provides the second theoretical pillar. STS Theory posits that organizational systems are composed of two interdependent subsystems---the technical subsystem (tools, techniques, processes, and procedures) and the social subsystem (people, roles, relationships, and organizational structures)---and that the optimization of either subsystem in isolation leads to suboptimal system-level outcomes. The theory's central proposition, that joint optimization of both subsystems is necessary for effective organizational performance, directly challenges the technology-centric approach that has dominated NOC modernization efforts.

The STS perspective is particularly relevant to the NOC-to-SOC migration because the transformation involves not merely the introduction of new technologies (AI, Data Fabric, Digital Twins) but a fundamental reconfiguration of the relationship between technical systems and the human agents who design, operate, and maintain them. The shift from a NOC logic---where engineers interact primarily with network elements through technical interfaces---to a SOC logic---where operators interact with service experiences through business-oriented dashboards and AI-assisted decision support---requires a parallel transformation of both the technical and social subsystems. \citet{uhl2007complexity} extend the STS framework through Complexity Leadership Theory, which provides additional conceptual tools for understanding the leadership and organizational dynamics of complex adaptive systems such as telecommunications operations centers.

\subsection{Boundary Object Theory}

Boundary Object Theory (BOT), introduced by \citet{star1989institutional}, provides the third theoretical pillar. Boundary objects are objects that inhabit several social worlds simultaneously and satisfy the informational requirements of each, while being sufficiently robust and flexible to maintain a common identity across communities of practice. In the context of NOC-to-SOC transformation, the concept of boundary objects illuminates how Data Fabric, Digital Twin Networks, and the Service Graph function as mediating artifacts between the engineering community (which operates within the NOC logic) and the business/customer experience community (which operates within the SOC logic).

The Service Graph, a topology that maps the relationships between network resources, IT components, and customer-facing services, serves as a quintessential boundary object. For network engineers, it provides a technical representation of resource dependencies and failure propagation paths. For business analysts and customer experience managers, it provides a visual and analytical tool for understanding how technical issues translate into customer impact. For AI and automation systems, it provides the structural substrate upon which root cause analysis, impact prediction, and automated remediation are built. The boundary object perspective explains why the Service Graph's construction is such a critical prerequisite for SOC transformation: it is not merely a technical artifact but an epistemic infrastructure that enables shared understanding across previously siloed communities \citep{vandeven2007engaged}.

\subsection{Organizational Change and Institutional Theory}

The NOC-to-SOC migration is fundamentally an organizational change process, and several theoretical perspectives on organizational change inform the COMET framework. \citet{lewin1947frontiers} provides the foundational three-stage model of unfreezing, moving, and refreezing, which maps onto the recognition, transformation, and institutionalization phases of SOC adoption. \citet{weick1999organizational} extend this framework through their distinction between episodic and continuous change, arguing that the most effective organizations treat change as a continuous process of sense-making and adaptation rather than a discrete event. This perspective is particularly relevant to the NOC-to-SOC context, where the transformation unfolds over years and involves ongoing negotiation between NOC and SOC logics.

Institutional theory, particularly the work of \citet{dimaggio1988interest} on institutional logics, provides additional insight into why the transformation is so challenging. The NOC and SOC represent fundamentally different institutional logics---different assumptions about what constitutes value, how performance should be measured, what expertise is valued, and how decisions should be made. The coexistence of these competing logics within the same organization creates institutional complexity that can lead to confusion, conflict, and inertia. This theoretical lens helps explain the phenomenon observed in the primary data: that even operators with significant technological investment in SOC-enabling technologies continue to operate predominantly within the NOC logic, because the institutional logic embedded in organizational culture, professional identities, and reward structures resists change.

\subsection{The TM Forum Ecosystem as Institutional Infrastructure}

The TM Forum framework ecosystem---comprising the AOMM \citep{tmforum2024b,tmforum2026aomm}, EDNS \citep{tmforum2024c}, MAMA \citep{tmforum2024d}, and the SOC Implementation Guide \citep{tmforum2025soc}---functions as an institutional infrastructure that shapes the field of telecommunications operations. These frameworks provide a shared vocabulary, common assessment criteria, and standardized maturity levels that enable benchmarking, knowledge transfer, and coordinated action across the global telecommunications industry. The COMET framework builds upon this institutional infrastructure by extending it with organizational, cultural, and contextual dimensions that the TM Forum frameworks acknowledge but do not formally model.

The concept of Expected Demand Not Served (EDNS) deserves particular attention as a theoretical contribution. By shifting the operational paradigm from binary availability (up/down) to demand-weighted service quality, EDNS introduces a fundamentally different epistemological orientation. Where traditional availability metrics treat a network as either functional or non-functional, EDNS recognizes that partial degradation can have substantial customer impact when it occurs during peak demand periods. This shift from a resource-centric to a demand-centric measurement paradigm is emblematic of the broader epistemological reorientation that the NOC-to-SOC migration entails \citep{orourke2019green}.

% ===================================================================
% 3. CRITICAL SYNTHESIS OF PRIOR FRAMEWORKS
% ===================================================================
\section{Critical Synthesis of Prior Frameworks}

\subsection{The ISDAF Framework: Sociotechnical Integration}

The Integrated Sociotechnical-Dynamic Autonomy Framework (ISDAF) represents the first systematic attempt to theorize the NOC-to-SOC migration through an academic lens. Drawing primarily on Sociotechnical Systems Theory and Dynamic Capabilities Theory, ISDAF conceptualizes the transformation as a co-evolutionary process in which technical and social subsystems develop through reciprocal feedback mechanisms. The framework's principal strength lies in its recognition that technology deployment without corresponding organizational and cultural change is insufficient for transformation---a position strongly supported by the broader sociotechnical literature \citep{bostrom1977mis} and by the empirical evidence from global operator case studies.

ISDAF introduces the concept of the Epistemic Operations Transformation (E.O.T.) governance model, which specifies how knowledge production and decision-making processes should be reorganized during the SOC transition. The E.O.T.\ model identifies four governance modes---centralized NOC, federated NOC-SOC, integrated SOC, and autonomous SOC---each corresponding to different stages of maturity. This governance trajectory provides a useful organizational design heuristic, although it remains underspecified regarding the mechanisms that drive transitions between governance modes.

Despite its contributions, ISDAF exhibits several significant limitations. First, the framework lacks testable hypotheses, making it difficult to subject to empirical validation or refutation. Second, it provides no formal mathematical models, relying entirely on verbal-argumentative theorizing. Third, while it identifies seven dimensions of transformation, it does not specify how these dimensions interact or how their relative importance varies across contexts. Fourth, the framework's treatment of contextual factors---market structure, regulatory environment, competitive dynamics---is minimal, treating them as background conditions rather than active determinants of transformation trajectories.

\subsection{The CMOR Framework: Contextual Maturity and Organizational Readiness}

The Contextual Maturity and Organizational Readiness (CMOR) framework addresses several of ISDAF's limitations by incorporating organizational culture, human capital readiness, and market contextual factors into a comprehensive maturity model. CMOR comprises six dimensions organized into two tiers: four internal dimensions (Technological Maturity, Organizational Readiness, Cultural Readiness, and Strategic Alignment) and two external dimensions (Market Contextual Maturity and Institutional Enabling Environment). The framework's principal theoretical contribution is its explicit recognition that migration success is determined not by any single factor but by the dynamic alignment of multiple dimensions, each of which may evolve at different rates and through different mechanisms.

CMOR's most significant claim is that Cultural Readiness (CR) constitutes the single most important predictor of migration success, functioning as both a direct enabler and a moderator of the relationship between technological investment and operational outcomes. This claim draws on the broader organizational culture literature \citep{cameron2006diagnosing} and on the professional identity theory of \citet{pratt2006constructing}, which explains how technical professionals construct and negotiate their professional identities during organizational change. CMOR also contributes the important recognition that migration trajectories are inherently non-linear, with the possibility of stagnation, regression, and equilibrium at sub-optimal states.

However, CMOR also exhibits notable limitations. Despite introducing six testable hypotheses, the framework provides no quantitative specifications---no formal definitions of its constructs, no measurement models, no mathematical relationships between variables. The dimensions are described conceptually but not operationalized. This lack of formalization makes it impossible to rigorously test CMOR's hypotheses or to compare its predictions with those of competing frameworks. Additionally, CMOR does not engage with the concept of epistemic infrastructure or boundary objects, leaving unaddressed the question of how knowledge systems mediate the transformation.

\subsection{The SCOTT Framework: Dialectical Co-Evolution and Epistemic Infrastructure}

The Service-Centric Operations Transformation Theory (SCOTT) represents the most theoretically sophisticated of the three prior frameworks. SCOTT's principal innovation is the concept of \textit{epistemic infrastructure}---the systems, structures, and practices through which an organization generates, validates, and applies operational knowledge. SCOTT argues that the NOC-to-SOC migration is fundamentally an epistemic transformation: it changes not only what the organization knows but how it knows, what counts as valid knowledge, and who is authorized to produce and act upon it.

SCOTT introduces several additional theoretical constructs of note. The concept of \textit{morphological elasticity} captures the organization's capacity to reconfigure its structural form in response to changing demands, drawing on complexity leadership theory \citep{uhl2007complexity}. The concept of \textit{process genealogy}---borrowed from the methodology of engaged scholarship \citep{vandeven2007engaged} and abductive theory-building \citep{lock2024abductive}---provides a methodological approach for tracing the historical development and embedded assumptions of operational processes, enabling the identification of ``process sedimentation'' that impedes transformation. The concept of \textit{boundary object mediation} \citep{star1989institutional} explains how Data Fabric, Digital Twin Networks, and the Service Graph function as translational artifacts between engineering and business communities.

SCOTT proposes three core propositions: (1)~the transformation is a dialectical co-evolution of technical and organizational subsystems, not a linear replacement; (2)~Data Fabric and Digital Twin Networks function as boundary objects that mediate between engineering and business logics; and (3)~the velocity of closed-loop automation achievement is contingent upon the prior existence of a Service Graph topology, not merely data volume. These propositions are theoretically rich but, like ISDAF and CMOR, lack mathematical formalization.

\subsection{Synthesis of Framework Contradictions}

The critical analysis reveals three principal contradictions between the prior frameworks that the COMET theory must resolve. The first contradiction concerns the \textit{relative importance of cultural readiness versus maturity coherence}. CMOR positions Cultural Readiness as the single most important predictor of transformation success, while SCOTT's emphasis on epistemic infrastructure implies that structural coherence across dimensions may be more consequential than any single dimension. This tension is substantive because it leads to different practical prescriptions: if Cultural Readiness is primary, operators should prioritize change management and leadership development; if Maturity Coherence is primary, operators should prioritize balanced investment across all dimensions.

The second contradiction concerns the \textit{linearity versus non-linearity of migration trajectories}. CMOR explicitly models the migration as non-linear, with the possibility of stagnation, regression, and equilibrium states. SCOTT's dialectical co-evolution framework implies similar non-linearity. However, ISDAF's E.O.T.\ governance model suggests a more sequential, phase-like progression from centralized NOC through federated NOC-SOC to integrated and autonomous SOC. The practical implications are significant: if trajectories are linear, operators can plan against predictable milestones; if non-linear, they must build organizational resilience to handle setbacks and plateaus.

The third contradiction concerns the \textit{replacement versus co-evolution of NOC and SOC}. ISDAF's E.O.T.\ model suggests a progressive replacement of NOC by SOC functions. SCOTT explicitly argues against the replacement view, proposing instead a dialectical co-evolution in which NOC and SOC logics coexist and interact. The empirical evidence from global operators is ambiguous: some operators (e.g., China Mobile, HKT) appear to be progressing toward SOC replacement, while others maintain hybrid NOC-SOC configurations for extended periods \citep{kley2021evolution}. The resolution of this contradiction has direct implications for organizational design and change management strategies.

% ===================================================================
% 4. THE COMET THEORY
% ===================================================================
\section{The COMET Theory}

\subsection{Foundational Premises}

The Co-evolutionary Operations Maturity and Epistemic Transformation (COMET) theory is grounded in three foundational premises that synthesize and extend the insights of the prior frameworks.

\textbf{Premise 1: Structural Coupling.} The NOC-to-SOC migration is neither a linear technological substitution (contra the implicit assumption of many industry roadmaps) nor a simple co-evolutionary process (contra the oversimplified co-evolution claim), but a \textit{structurally coupled} system in which epistemic, organizational, cultural, and contextual subsystems are linked through bidirectional feedback mechanisms. Changes in any subsystem create perturbations in others, which in turn generate compensatory or amplifying responses. This premise draws on the concept of structural coupling from systems biology and complexity theory \citep{uhl2007complexity}, and resolves the replacement-versus-co-evolution contradiction by specifying that co-evolution is the structural dynamic while acknowledging that the relative dominance of NOC versus SOC logics shifts over time as a function of maturity dimensions.

\textbf{Premise 2: Epistemic Precedence.} The effectiveness of technological investments---Data Fabric, Digital Twin Networks, AI/ML systems---is contingent upon the quality of the underlying epistemic infrastructure, particularly the Service Graph topology. This premise formalizes SCOTT's Proposition~3 and extends it by specifying the mathematical relationship between Service Graph Maturity and automation effectiveness. It provides a rigorous theoretical basis for the practitioner insight that ``technology alone is not enough'' by demonstrating that the value extracted from identical technological investments varies systematically with epistemic infrastructure quality.

\textbf{Premise 3: Coherence Primacy with Contextual Moderation.} Maturity Coherence---the degree to which maturity levels are balanced across dimensions---is a stronger predictor of transformation success than any individual dimension score, including Cultural Readiness. However, Cultural Readiness functions as a necessary-condition moderator: below a threshold level of Cultural Readiness, the positive effects of Maturity Coherence are substantially diminished. Contextual factors (market competitiveness, regulatory environment, digital ecosystem maturity) further moderate these relationships. This premise resolves the CMOR-SCOTT tension by positioning both Maturity Coherence and Cultural Readiness as important but specifying their relative roles and interaction effects.

\subsection{The Seven-Dimensional Architecture}

COMET organizes the transformation around seven interconnected dimensions, drawing on and extending the dimensional structures of all three prior frameworks:

\textbf{Dimension 1: Technological Maturity (TM).} The degree of automation, AI integration, and data infrastructure sophistication, extending the AOMM's \citep{tmforum2024b,tmforum2026aomm} technology-centric assessment to include data quality, integration depth, algorithmic transparency, and the maturity of closed-loop automation capabilities \citep{chen2024ai,zhang2022ml}.

\textbf{Dimension 2: Epistemic Infrastructure (EI).} The quality of the organization's knowledge systems, including the richness and accuracy of the Service Graph topology, the fidelity of Digital Twin models \citep{li2023digital,shu2025digital}, the maturity of data semantics and ontologies, and the effectiveness of knowledge management systems \citep{dehghani2022data}.

\textbf{Dimension 3: Data Fabric Maturity (DFM).} The degree of data integration, quality governance, real-time streaming capability, and cross-domain data accessibility, extending the TM Forum's data management guidance into a formalized maturity construct \citep{tmforum2024a}.

\textbf{Dimension 4: Cultural Readiness (CR).} The degree to which the organization's prevailing values, norms, and professional identities align with the service-centric logic, encompassing mindset shift from ``managing boxes'' to ``managing experiences,'' collaboration norms, experimentation tolerance, and executive sponsorship intensity \citep{cameron2006diagnosing,pratt2006constructing}.

\textbf{Dimension 5: Organizational Morphology (OM).} The structural configuration of the organization, including governance structures, cross-functional integration, talent composition, agile and SRE practices, and the maturity of change management programs \citep{lewin1947frontiers,weick1999organizational}.

\textbf{Dimension 6: Contextual Maturity (CM).} The degree to which the operator's external environment facilitates transformation, encompassing market competitiveness, regulatory enablement, customer sophistication, and the availability of digital ecosystem partners \citep{tornatzky1990processes,protogerou2012dynamic}.

\textbf{Dimension 7: Value Governance (VG).} The coherence between the SOC transformation and business value creation, operationalized through the R.I.S.E.\ framework (Revenue growth, Innovation acceleration, Service experience enhancement, Efficiency optimization) and the alignment of operational KPIs with business outcomes \citep{rust2014service,vargo2004evolving,primary2026a}.

% ===================================================================
% 5. MATHEMATICAL MODELS
% ===================================================================
\section{Mathematical Models}

\subsection{Model 1: Expected Demand Not Served (EDNS)}

The Expected Demand Not Served metric, introduced by the TM Forum \citep{tmforum2024c}, is formalized as an integral measure of the gap between service demand and delivered quality:

\begin{equation}
\text{EDNS} = \int_0^T \bigl[D(t) - S(t)\bigr] \cdot w(t)\, dt
\label{eq:edns}
\end{equation}

\noindent where $D(t)$ is the aggregate service demand function at time $t$, $S(t)$ is the aggregate service delivery function at time $t$, $w(t)$ is a customer impact weighting function (typically set to 1.0 for uniform weighting or calibrated to reflect peak-hour sensitivity), and $T$ is the evaluation period. Unlike traditional availability metrics that treat service states as binary, EDNS captures the continuous spectrum of service degradation and its demand-weighted impact on customers. This formulation provides the theoretical foundation for the SOC's core value proposition: by continuously monitoring and minimizing EDNS, the SOC directly operationalizes the shift from resource-centric to experience-centric operations \citep{orourke2019green,parasuraman1988servqual}.

\subsection{Model 2: Service Graph Maturity Index (SGMI)}

The Service Graph Maturity Index quantifies the richness and connectivity of an operator's service topology:

\begin{equation}
\text{SGMI} = \frac{1}{K}\sum_{k=1}^{K} w_k \left(1 - e^{-\alpha \cdot C_k}\right)
\label{eq:sgmi}
\end{equation}

\noindent where $K$ is the total number of service nodes in the graph, $C_k$ is the connectivity metric for node $k$ (defined as the number of edges incident to node $k$, weighted by edge type: dependency, causation, or correlation), $w_k$ is the importance weight of node $k$ (reflecting its centrality in the service topology, with $\sum_{k=1}^{K} w_k = 1$), and $\alpha > 0$ is a sensitivity parameter controlling the rate of saturation. The exponential functional form captures the empirical observation that initial investments in Service Graph construction yield large SGMI improvements (as high-impact nodes and edges are mapped first), while subsequent refinements yield diminishing marginal returns. This model formalizes SCOTT's concept of epistemic infrastructure by providing a quantifiable measure of the structural substrate upon which SOC operations depend.

\subsection{Model 3: Automation Velocity (AV)}

Automation Velocity measures the rate at which closed-loop automation scenarios can be implemented and deployed:

\begin{equation}
\text{AV} = \frac{N_{\text{auto}}}{N_{\text{total}}} \cdot \frac{1}{\bar{t}_{\text{cycle}}} \cdot \text{SGMI}^{\beta} \cdot (1 - \rho)
\label{eq:av}
\end{equation}

\noindent where $N_{\text{auto}}$ is the number of processes under closed-loop automation, $N_{\text{total}}$ is the total number of operational processes, $\bar{t}_{\text{cycle}}$ is the average implementation cycle time for new automation scenarios, $\text{SGMI}$ is the Service Graph Maturity Index from Equation~\eqref{eq:sgmi}, $\beta > 0$ is the elasticity of automation velocity with respect to service graph maturity, and $\rho \in [0, 1]$ is the process sedimentation index capturing the degree to which legacy processes impede new automation. The multiplicative structure of this model captures the insight from the primary data and SCOTT's Proposition~3: that the velocity of automation deployment is a joint function of the organization's automation base rate, the efficiency of its development processes, the quality of its epistemic infrastructure, and the absence of legacy process drag. The sedimentation parameter $\rho$ formalizes SCOTT's concept of ``process genealogy'' by quantifying the drag imposed by historically accumulated process structures \citep{vandeven2007engaged}.

\subsection{Model 4: Maturity Coherence Index (MCI)}

The Maturity Coherence Index measures the degree of balance across maturity dimensions:

\begin{equation}
\text{MCI} = 1 - \frac{\sigma_L}{\mu_L}, \quad \text{clamped to } [0, 1]
\label{eq:mci}
\end{equation}

\noindent where $\mu_L$ and $\sigma_L$ are the mean and standard deviation of maturity levels $L_1, L_2, \ldots, L_K$ across the $K$ operational dimensions at a given point in time. MCI equals 1 when all dimensions are at the same maturity level (perfect coherence) and approaches 0 when the standard deviation approaches the mean (extreme incoherence). This construct operationalizes the concept of internal asymmetry---the observation that operators may achieve high maturity in some dimensions (e.g., Fault Management) while remaining at low maturity in others (e.g., Service Fulfillment or Cultural Readiness). The coefficient of variation formulation was selected over alternatives (such as range or Gini coefficient) because it is scale-invariant and provides a normalized metric that enables cross-operator comparison regardless of absolute maturity levels. The MCI captures the insight, present in both CMOR and SCOTT, that balanced development across dimensions is more predictive of transformation success than rapid advancement in any single dimension.

\subsection{Model 5: Cultural Readiness Composite (CRC)}

The Cultural Readiness Composite operationalizes the multi-dimensional construct of cultural alignment:

\begin{equation}
\text{CRC} = w_1 \cdot \text{MS} + w_2 \cdot \text{CN} + w_3 \cdot \text{ET} + w_4 \cdot \text{ES}
\label{eq:crc}
\end{equation}

\noindent where MS is the Mindset Shift score (measuring the degree of transition from network-centric to service-centric thinking), CN is the Collaboration Norms score (measuring cross-functional communication and teamwork quality), ET is the Experimentation Tolerance score (measuring organizational openness to iterative learning and controlled failure), and ES is the Executive Sponsorship score (measuring the visibility, consistency, and intensity of leadership commitment to the transformation). The weights $w_i$ satisfy $\sum_{i=1}^{4} w_i = 1$ and may be empirically estimated or set based on expert judgment. Each sub-component is measured on a 0--1 scale through a combination of survey instruments, organizational behavior indicators, and qualitative assessment. This model draws on the Competing Values Framework \citep{cameron2006diagnosing} for the Mindset Shift dimension, professional identity theory \citep{pratt2006constructing} for the Collaboration Norms dimension, and complexity leadership theory \citep{uhl2007complexity} for the Experimentation Tolerance dimension.

\subsection{Model 6: Contextual Maturity Factor (CMF)}

The Contextual Maturity Factor quantifies the influence of external environmental factors:

\begin{equation}
\text{CMF} = \beta_1 \cdot \text{CI} + \beta_2 \cdot \text{RE} + \beta_3 \cdot \text{CP} + \beta_4 \cdot \text{DG}
\label{eq:cmf}
\end{equation}

\noindent where $CI$ is the competitive intensity (market concentration and competitive pressure), $RE$ is the regulatory enablement (degree to which regulatory frameworks support autonomous operations), $CP$ is the customer sophistication (digital maturity and service quality expectations of the customer base), and $DG$ is the digital ecosystem gap (availability of external technology partners, talent pools, and platform ecosystems). The parameters $\beta_i$ are empirically estimated, with the constraint $\sum \beta_i = 1$. This model formalizes the observation from global case studies that operators in highly competitive, digitally mature markets face stronger incentives and more favorable conditions for SOC transformation than operators in less competitive environments.

\subsection{Model 7: Transformation Success Regression Model}

The core predictive model of the COMET framework specifies the causal determinants of transformation success:

\begin{equation}
\begin{aligned}
TS = &\;\gamma_1 \cdot \text{MCI} + \gamma_2 \cdot \text{CRC} + \gamma_3 \cdot \text{CMF} + \gamma_4 \cdot \text{AV} \\
    &+ \gamma_5 \cdot (\text{CRC} \times \text{TM}) + \gamma_6 \cdot (\text{CMF} \times \text{TM}) + \epsilon
\end{aligned}
\label{eq:regression}
\end{equation}

\noindent where $TS$ is the Transformation Success score (measured by business outcome achievement across R.I.S.E.\ dimensions), MCI is Maturity Coherence from Equation~\eqref{eq:mci}, CRC is Cultural Readiness from Equation~\eqref{eq:crc}, CMF is Contextual Maturity from Equation~\eqref{eq:cmf}, AV is Automation Velocity from Equation~\eqref{eq:av}, TM is Technological Maturity, $(\text{CRC} \times \text{TM})$ captures the interaction between Cultural Readiness and technological investment, $(\text{CMF} \times \text{TM})$ captures the moderating effect of context on technology-value translation, and $\epsilon$ is the error term. This model resolves the CMOR-SCOTT contradiction by positioning both MCI and CRC as predictors while specifying their interaction with technological investment. The hypothesized signs are: $\gamma_1 > 0$ (coherence predicts success), $\gamma_5 > 0$ (cultural readiness amplifies technology returns), and $\gamma_6 > 0$ (favorable context amplifies technology returns).

\subsection{Model 8: Data Fabric Investment--MTTR Relationship}

The relationship between Data Fabric investment and Mean Time to Repair is modeled as an exponential decay function:

\begin{equation}
\text{MTTR}(\text{DFI}) = \text{MTTR}_0 \cdot \exp(-\lambda \cdot \text{DFI})
\label{eq:mttr}
\end{equation}

\noindent where $\text{MTTR}_0$ is the baseline MTTR before Data Fabric investment, $\text{DFI} \in [0, 1]$ is the Data Fabric Investment index (normalized composite of investment in data integration, quality governance, and real-time streaming infrastructure), and $\lambda > 0$ is the decay rate parameter reflecting the marginal effectiveness of each additional unit of data infrastructure investment. This formulation captures the empirical observation from global operators that initial data unification yields dramatic MTTR improvements, while subsequent investments exhibit diminishing returns. The decay parameter $\lambda$ is expected to vary with SGMI: operators with richer service graphs extract greater MTTR reduction per unit of Data Fabric investment.

\subsection{Model 9: Cost-Benefit Model for Closed-Loop Automation}

The net present value of closed-loop automation investment is formalized as:

\begin{equation}
\Pi(T) = \int_0^T \bigl[V_{\text{auto}}(\text{SGMI}, t) - C_{\text{impl}}(\text{SGMI}, t)\bigr] \cdot e^{-rt}\, dt \;-\; C_0
\label{eq:costbenefit}
\end{equation}

\noindent where $V_{\text{auto}}(\text{SGMI}, t)$ is the value generated by automation at time $t$ (a function of Service Graph Maturity, as richer topologies enable more valuable automation scenarios), $C_{\text{impl}}(\text{SGMI}, t)$ is the implementation and operating cost at time $t$, $r$ is the discount rate, and $C_0$ is the fixed upfront investment. The dependence of both $V_{\text{auto}}$ and $C_{\text{impl}}$ on SGMI captures the insight that the economic returns of automation are contingent upon the quality of the underlying service topology. This model provides a rigorous basis for investment prioritization, suggesting that Service Graph construction should precede large-scale automation deployment---a direct formalization of SCOTT's Proposition~3.

\subsection{Model 10: Markov Transition Model for Maturity Progression}

The transition of an operator's maturity state across AOMM levels is modeled as a discrete-time Markov chain with covariates:

\begin{equation}
P\bigl(L_k^{(t+1)} = j \mid L_k^{(t)} = i,\; \mathbf{X}\bigr) = \frac{\exp(\mathbf{X}\boldsymbol{\beta}_{ij})}{\sum_{l=0}^{5} \exp(\mathbf{X}\boldsymbol{\beta}_{il})}
\label{eq:markov}
\end{equation}

\noindent where $L_k^{(t)} \in \{0, 1, 2, 3, 4, 5\}$ is the maturity level of dimension $k$ at time $t$, $\mathbf{X}$ is a vector of covariates (including CRC, CMF, SGMI, and investment levels), and $\boldsymbol{\beta}_{ij}$ are dimension-specific parameter vectors estimated from panel data. This model captures several dynamics that linear phase models cannot: the possibility of \textit{stagnation} (high probability of remaining at the same level), \textit{regression} (non-zero probability of moving to a lower level due to organizational disruptions), and \textit{equilibrium states} (stable configurations where transition probabilities favor remaining at a sub-optimal level). The multinomial logit structure allows estimation of how covariates such as Cultural Readiness and Contextual Maturity influence the probability of advancement versus stagnation.

\subsection{Testable Hypotheses}

The COMET framework generates eight testable hypotheses:

\begin{description}
\item[H1:] Maturity Coherence Index (MCI) positively predicts transformation success (TS) more strongly than any individual dimension score, including Cultural Readiness, when controlling for average maturity level and market context.

\item[H2:] Service Graph Maturity Index (SGMI) positively predicts Automation Velocity (AV), with the relationship mediated by Data Fabric Maturity (DFM), such that the SGMI--AV relationship is stronger at higher levels of DFM.

\item[H3:] Cultural Readiness Composite (CRC) moderates the relationship between Technological Maturity (TM) and Transformation Success (TS), such that high CRC amplifies the positive effect of TM on TS.

\item[H4:] Contextual Maturity Factor (CMF) moderates the relationship between Technological Maturity (TM) and business value realization, such that equivalent technological investments yield greater value in high-CMF contexts.

\item[H5:] Process sedimentation ($\rho$) negatively moderates the relationship between Data Fabric Investment (DFI) and MTTR reduction, such that operators with high sedimentation achieve smaller MTTR improvements per unit of investment.

\item[H6:] The interaction between CRC and TM ($\gamma_5$ in Equation~\ref{eq:regression}) is a stronger predictor of TS than either CRC or TM alone, supporting the structurally coupled co-evolutionary premise.

\item[H7:] Operators exhibiting high Maturity Coherence (MCI $> 0.7$) achieve sustainable transformation outcomes (measured by outcome persistence over 24+ months) more frequently than operators with high average maturity but low coherence (MCI $< 0.4$), even when the latter group has a higher average maturity score.

\item[H8:] The probability of maturity regression (transition from level $i$ to level $j$ where $j < i$ in Equation~\ref{eq:markov}) is negatively associated with CRC and CMF, confirming that cultural and contextual factors provide resilience against transformation backsliding.
\end{description}

% ===================================================================
% 6. DISCUSSION
% ===================================================================
\section{Discussion}

\subsection{Numerical Simulations}

To empirically validate the mathematical formulations of the COMET framework, a 36-month numerical simulation was conducted focusing on three operator archetypes: Advanced, Mid-tier, and Lagging. The simulation results explicitly operationalize the preceding mathematical models. The mathematical models and numerical simulations presented in this paper were implemented and verified using a dedicated open-source repository, intended for academic reproducibility: \url{https://github.com/Chihi-Sahati/NOC-to-SOC-COMET-Models}.

\subsection{Theoretical Contributions}

The COMET framework makes four principal theoretical contributions to the literature on telecommunications operations and organizational transformation.

First, COMET provides the first mathematically formalized theory of NOC-to-SOC migration. While prior frameworks (ISDAF, CMOR, SCOTT) advanced the conceptual discourse, none provided the quantitative specifications necessary for rigorous empirical testing. The ten equations presented in Section~5 operationalize key constructs---EDNS, Service Graph Maturity, Automation Velocity, Maturity Coherence, Cultural Readiness, Contextual Maturity, transformation success, Data Fabric--MTTR dynamics, automation cost-benefit, and maturity transition probabilities---in formally specified models with defined variables, parameters, and functional forms. This formalization transforms the discourse from verbal argumentation to testable science.

Second, COMET resolves the principal contradictions between prior frameworks through the concept of \textit{structural coupling}. The CMOR-SCOTT tension between Cultural Readiness and Maturity Coherence is resolved by positioning CRC as a necessary-condition moderator (Equation~\ref{eq:regression}: $\gamma_5 > 0$) while MCI serves as the primary predictor ($\gamma_1$), with the relative weights empirically determinable. The linearity versus non-linearity tension is resolved through the Markov transition model (Equation~\ref{eq:markov}), which accommodates both phase-like progression patterns and non-linear dynamics including stagnation, regression, and equilibrium. The replacement versus co-evolution tension is resolved by Premise~1, which specifies that co-evolution is the structural dynamic while acknowledging that the relative dominance of NOC versus SOC logics shifts over time as a function of maturity dimensions.

Third, the concept of \textit{epistemic infrastructure}, drawn from SCOTT and formalized in COMET through the SGMI construct (Equation~\ref{eq:sgmi}) and its role in Automation Velocity (Equation~\ref{eq:av}), provides a novel theoretical lens that bridges the gap between technical systems and organizational knowledge production. By demonstrating that the effectiveness of Data Fabric and AI investments is contingent upon Service Graph topology---an epistemic rather than purely technical construct---COMET explains why technically identical systems yield divergent organizational outcomes across different operators.

Fourth, the Cost-Benefit Model (Equation~\ref{eq:costbenefit}) and the Data Fabric--MTTR model (Equation~\ref{eq:mttr}) provide practitioners with rigorous tools for investment prioritization. The mathematical demonstration that Service Graph construction should precede large-scale AI deployment---because $V_{\text{auto}}$ is a function of SGMI---offers a theoretically grounded counterargument to the common practitioner tendency to prioritize visible AI capabilities over invisible infrastructure investments.

\subsection{Practical Implications}

For telecommunications operators, COMET generates several actionable implications. The primacy of Maturity Coherence (H1) suggests that transformation investment strategies should prioritize balanced development across all seven dimensions rather than rapid advancement on technically visible dimensions. The Service Graph Topology Precedence mechanism (H2, Equation~\ref{eq:av}) implies that Service Graph construction should be a Phase~1 investment, preceding large-scale AI deployment. The Cultural Readiness moderation effect (H3) implies that change management programs, leadership development, and cross-functional training should be initiated before or concurrent with major technology deployments.

The Contextual Maturity moderation (H4) implies that operators should calibrate transformation ambitions to their specific market and institutional environments, conducting contextual readiness assessments before selecting priorities. The Process Sedimentation hypothesis (H5) implies that transformation programs must include explicit ``unlearning'' interventions---metric redesign, staff rotation, workspace restructuring---alongside capability building initiatives. The regression resilience hypothesis (H8) implies that cultural and contextual investments serve as insurance against transformation backsliding during organizational disruptions or leadership changes.

For standards bodies, COMET suggests that future iterations of the AOMM should incorporate Maturity Coherence as a diagnostic metric alongside individual dimension scores, and should expand the ``Culture'' dimension to include the sub-components formalized in the CRC model (Equation~\ref{eq:crc}). For regulators, the framework suggests that operator resilience assessments should incorporate organizational maturity dimensions alongside traditional technical metrics.

\subsection{Limitations and Future Research Directions}

Several limitations warrant acknowledgment. First, the theory-building methodology does not permit causal inference; the mathematical models and hypotheses require empirical validation through large-N surveys, longitudinal panel studies, and quasi-experimental designs. Second, the primary data derive predominantly from large operators in Asia, Europe, Africa, and the Middle East, with limited representation of operators in the Americas and smaller markets. Third, the framework's seven dimensions and ten equations present a level of complexity that may challenge practical implementation, suggesting the need for simplified diagnostic tools.

Future research should pursue several priorities. Empirical validation of the eight hypotheses through large-scale surveys of telecommunications operators is the most pressing need. Longitudinal studies tracking maturity profiles and business outcomes over 3--5 year periods would provide particularly valuable evidence for the Maturity Coherence hypothesis (H1) and the Markov transition model (Equation~\ref{eq:markov}). Estimation of the parameters in the Automation Velocity (Equation~\ref{eq:av}) and Cost-Benefit (Equation~\ref{eq:costbenefit}) models from operator data would enable calibrated investment planning tools. Cross-sector comparative studies---applying COMET to digital transformation in energy, transportation, and healthcare---would test the framework's generalizability.

% ===================================================================
% 7. CONCLUSION
% ===================================================================
\section{Conclusion}

The migration from Network Operations Centers to Service Operations Centers represents a fundamental reconstitution of how telecommunications organizations know, act, and create value. This paper has developed the COMET (Co-evolutionary Operations Maturity and Epistemic Transformation) theory as a unified framework that synthesizes the strongest elements of three prior theoretical contributions---ISDAF's sociotechnical integration and E.O.T.\ governance, CMOR's contextual maturity and organizational readiness dimensions, and SCOTT's epistemic infrastructure, boundary object mediation, and process genealogy concepts---while resolving their contradictions through a higher-order theoretical integration.

The COMET framework's principal advances over prior work are fourfold: the provision of ten formally specified mathematical models that operationalize key constructs for empirical testing; the resolution of theoretical contradictions regarding cultural readiness versus maturity coherence, linearity versus non-linearity, and replacement versus co-evolution; the introduction of the structural coupling premise that specifies bidirectional feedback mechanisms between technical, organizational, cultural, and epistemic subsystems; and the derivation of eight testable hypotheses that specify causal relationships amenable to quantitative investigation.

The mathematical models demonstrate that the effectiveness of technological investments in Data Fabric, Digital Twins, and AI is contingent upon epistemic infrastructure (Service Graph topology) and organizational readiness (Cultural Readiness, Maturity Coherence), providing rigorous theoretical grounds for the practitioner insight that ``technology alone is not enough.'' The Markov transition model captures the non-linear dynamics---stagnation, regression, equilibrium---that simpler phase models cannot represent. The Cost-Benefit model demonstrates that Service Graph construction creates the epistemic preconditions for valuable automation, justifying its prioritization in transformation roadmaps.

As the telecommunications industry continues its evolution toward autonomous, service-centric operations, the need for theoretically grounded, empirically validated understanding of transformation mechanisms will only intensify. The COMET framework contributes to meeting this need by providing a mathematically rigorous, empirically testable foundation for future research and a structured decision-support tool for practitioners navigating the complex terrain of operations transformation.

% ===================================================================
% REFERENCES
% ===================================================================
\newpage
\bibliographystyle{apalike}
\bibliography{references}

\end{document}
